\chapter{WStarAlgebras and their Topologies}

\section{Basic Definitions and Results}

In this section we collect basic definitions and results about $C^{*}$--algebras and $W^{*}$--algebras.

An associative algebra $\mathcal{A}$ over the complex numbers is called a \textbf{normed algebra}
  if there is a map assigning for every $x\in \mathcal{A}$ a real number $\|x\|$, and this map satisfies:
  \begin{enumerate}
  \item $\|x\|\ge 0$ and $\|x\|=0 \iff x=0$;
  \item $\|x+y\|\le \|x\|+\|y\|$;
  \item $\|\lambda x\|= |\lambda|\|x\|$;
  \item $\|xy\|\le \|x\|\|y\|$.
  \end{enumerate}
  If $\mathcal{A}$ is also complete with respect to the norm, it is called a \textbf{Banach algebra}.
  A mapping $*: \mathcal{A}\to \mathcal{A}$ is called an \textbf{involution} if it satisfies:
  \begin{enumerate}
  \item $(x^{*})^{*}=x$;
  \item $(x+y)^{*}=x^{*}+y^{*}$;
  \item $(xy)^{*}=y^{*}x^{*}$;
  \item $(\lambda x)^{*}=\overline{\lambda}x^{*}$.
  \end{enumerate}
  where $\lambda$ above is a complex number.

\begin{definition}
  \label{def:cstar_def}
  \lean{CStarAlgebra,NonUnitalCStarAlgebra}
  \mathlibok
  A Banach algebra $\mathcal{A}$ with an involution is a $\mathbf{C^{*}}$\textbf{--algebra} if the relation
  $\|x^{*}x\|=\|x\|^{2}$ holds for every $x\in \mathcal{A}$. A $C^{*}$--algebra is \textbf{unital} if it has a multiplicative
  identity.
\end{definition}

\begin{lemma}
  \label{lem:star_pres_norm_Sak_1_1_6} (Sak 1.1.6)
  \lean{CStarRing.to_normedStarGroup}
  \mathlibok
  \uses{def:cstar_def}
  The involution on a $C^{*}$--algebra preserves the norm.
\end{lemma}

\begin{proof}
  \leanok
  Note $\|x\|^{2}=\|x^{*}x\|\le \|x^{*}\|\|x\|$, and $\|x^{*}\|^{2}\le \|xx^{*}\|\le \|x\|\|x^{*}\|$.
\end{proof}

\begin{theorem}
  \label{thm:comm_cstar_cts_fcns_Sak_1_2_1} (Sak 1.2.1)
  \lean{gelfandStarTransform}
  \mathlibok
  \uses{def:cstar_def}
  Every commutative, unital $C^{*}$--algebra $\mathcal{A}$ is $*$--isomorphic to $C(K)$,
  the algebra of complex-valued continuous functions on $K$, where $K$ is the compact Hausdorff
  space of maximal ideals of $\mathcal{A}$. (Note $C(K)$ is itself a $C^{*}$--algebra with pointwise
  complex conjugation as $*$ and the uniform norm.)
\end{theorem}

\begin{proof}
  \leanok
  The Gelfand transform sends an element to its evaluations at the characters of $\mathcal A$.
  The $C^{*}$--identity and the spectral-radius formula make this homomorphism isometric, while
  characters take real values on self-adjoint elements, so it preserves the involution. Its range
  is therefore a closed, self-adjoint, unital subalgebra of $C(K)$ which separates points.
  The Stone--Weierstrass theorem shows that this range is all of $C(K)$.
\end{proof}

\begin{corollary}
  \label{cor:spec_sa_homeo_spec_cstar}
  \lean{cfcₙHom,range_cfcₙHom,IsStarNormal.instNonUnitalIsometricContinuousFunctionalCalculus}
  \mathlibok
  Let $a$ be a self-adjoint element of a possibly non-unital $C^{*}$--algebra $\mathcal A$.
  The non-unital continuous functional calculus embeds the complex-valued continuous functions on
  the pointed quasispectrum $\sigma_{n}(a)$ which vanish at its distinguished point $0$ isometrically
  into $\mathcal A$, with image the elemental non-unital $C^{*}$--subalgebra generated by $a$.
  Consequently, if $a$ generates $\mathcal A$, this map identifies that function algebra with
  $\mathcal A$.
\end{corollary}

\begin{proof}
  \leanok
  Apply the complex non-unital continuous functional calculus to the normal element $a$. Its
  defining homomorphism is an isometry, and the range theorem identifies its image with the
  elemental non-unital star subalgebra generated by $a$. This range is all of $\mathcal A$ when
  $a$ generates $\mathcal A$.
\end{proof}

\begin{lemma}
  \label{lem:one_is_extreme}
  \lean{CStarAlgebra.one_mem_extremePoints_unitClosedBall}
  \leanok
  \uses{thm:comm_cstar_cts_fcns_Sak_1_2_1,lem:star_pres_norm_Sak_1_1_6}
  In a unital $C^{*}$--algebra $\mathcal{A}$, the identity $1$ is an extreme point of the closed unit ball $S$.
\end{lemma}

\begin{proof}
  \leanok
  Suppose $a,b\in S$ and $1=\frac{a+b}{2}$. Then, since $1=1^{*}$, we have $1=\frac{a^{*}+b^{*}}{2}$ and one checks that $1=\frac{c+d}{2}$,
  where $c=\frac{a+a^{*}}{2}$ and $d=\frac{b+b^{*}}{2}$. Now, $d=2-c$ and so $d$ and $c$ commute and each are self-adjoint (hence normal)
  and therefore $C^{*}\{1,d,c\}$ is commutative and can be identified with $C(K)$ for a compact Hausdorff space $K$.
  The self-adjointness of $c$ and $d$, the triangle inequality and Lemma \ref{lem:star_pres_norm_Sak_1_1_6} imply
  that the range of the associated functions $f,g$ are in $[-1,1]$. In order for $f(t)+g(t)=2$ for all $t\in K$
  it must be that $f(t)$ and $g(t)$ are always $1$, and thus $c=d=1$. Now $2=2c=a+a^{*}$ and so $a=2-a^{*}$ and so $a$ is normal and
  so the unital $C^{*}$--subalgebra generated by $a$ is commutative.
  Again represent the unital $C^{*}$--algebra as $C(K)$ for some compact Hausdorff $K$. Let $f$ be the
  function corresponding to $a$. We have $\Re(f)(t)=1$ and $|f(t)|\le 1$ for all $t$, which
  forces $\Im(f)(t)=0$. Therefore $a=1$.
\end{proof}

\begin{lemma}
  \label{lem:fake_approx_unit}
  \lean{CStarAlgebra.exists_seq_tendsto_mul_left_sq}
  \leanok
  Let $a$ be an element of a possibly non-unital $C^{*}$--algebra $\mathcal A$. There exists a
  sequence $(y_n)$ of positive contractions such that $\|ay_n-a\|\to 0$ and
  $\|ay_n^{2}-a\|\to 0$ as $n\to\infty$.
\end{lemma}

\begin{proof}
  \leanok
  Apply the increasing approximate-unit filter of $\mathcal A$. For each $n$,
  choose a positive contraction $y_n$ for which
  $\|ay_n-a\|<(n+1)^{-1}$. This gives the first convergence. Moreover,
  \[
    ay_n^2-a=(ay_n-a)y_n+(ay_n-a),
  \]
  so contractivity of $y_n$ bounds the second error by twice the first.
\end{proof}

\begin{lemma}
  \label{lem:extreme_partial_isom}
  \lean{isStarProjection_star_mul_self_of_mem_extremePoints_unitClosedBall}
  \leanok
  If $x$ is an extreme point of the norm unit ball of a $C^{*}$--algebra $\mathcal{A}$, then $x^{*}x$ is a projection. I.e. $x$ is a partial
  isometry.
\end{lemma}

\begin{proof}
  \leanok
\end{proof}


\begin{theorem}
  \label{thm:cstar_unital_iff_extreme_Sak_1_6_1} (Sakai 1.6.1)
  \lean{CStarAlgebra.isUnital_iff}
  \leanok
  \uses{lem:one_is_extreme, lem:extreme_partial_isom}
  A $C^{*}$--algebra is unital if and only if its closed unit ball has an extreme point.
\end{theorem}

\begin{proof}
  \leanok
\end{proof}

\begin{proposition}
  \label{prop:extr_pos_projs_Sak_1_6_2} (Sakai 1.6.2)
  \lean{isStarProjection_iff_mem_extremePoints_setOf_nonneg_inter_unitClosedBall}
  \leanok
  Let $P$ be the set of positive elements in a $C^{*}$--algebra $\mathcal{A}$ with closed unit ball $S$. Then the extreme
  points of $P\cap S$ are the projections of $\mathcal{A}$.
\end{proposition}

\begin{proof}
  \leanok
\end{proof}

\begin{proposition}
  \label{prop:extr_sa_sa_unitaries_Sak_1_6_3} (Sakai 1.6.3)
  \lean{mem_extremePoints_isSelfAdjoint_and_mem_unitClosedBall_iff_isSelfAdjoint_and_mem_unitary}
  \leanok
  \uses{prop:extr_pos_projs_Sak_1_6_2}
  Let $\mathcal{A}^{s}$ be the set of self-adjoint elements in a $C^{*}$--algebra $\mathcal{A}$ with closed unit ball $S$. Then
  the extreme points of $\mathcal{A}^{s}\cap S$ is the set of all self-adjoint unitary elements of $\mathcal{A}$.
\end{proposition}

\begin{proof}
  \leanok
\end{proof}

\begin{definition}
  \label{def:wstar_def}
  \uses{def:cstar_def}
  \lean{WStarAlgebra,Predual}
  \leanok
  A $C^{*}$--algebra $M$ is called a $\mathbf{W^{*}}$\textbf{--algebra} if it is a dual Banach space, i.e. if there exists a Banach space
  $M_{*}$ so that $(M_{*})^{*}$ is isometrically isomorphic to $M$. We call $M_{*}$ a \textbf{predual} of $M$.
\end{definition}

Note: It is premature to fix a notation like $M_{*}$ in the above, but the aim of this project is to prove that any two
preduals of $M$ are isomorphic as Banach spaces, so we take the liberty with Sakai to ``fix a reference predual'' at the outset.

\begin{lemma}
    \label{lem:sigma_closed_star_subalg_wstar} (Sakai 1.1.4)
    \lean{WStarAlgebra.of_isClosed_submodule,StarSubalgebra.IsUltraweakClosed.wStarAlgebra}
    \leanok
    \uses{def:wstar_def,def:sigma_top}
  An ultraweakly closed star subalgebra $N$ of a $W^{*}$--algebra $M$ is also a $W^{*}$--algebra, since $(M_{*}/N^{0})^{*}=N$, where
  $N^{0}$ is the polar (annihilator) of $N$ in $M_{*}$, i.e. the set of $\varphi\in M_{*}$ such that $\varphi(x)=0$ for all $x\in N$.
\end{lemma}

\begin{proof}
  \leanok
  The preannihilator $N^0$ is norm closed, and the bipolar theorem identifies $N$ with the
  annihilator of $N^0$. The continuous dual of the quotient $M_*/N^0$ is isometrically the
  annihilator of $N^0$, so an explicit isometry from the algebra to the ambient closed submodule
  supplies the required predual without relying on definitional equality of subtype structures.
\end{proof}

\section{The Ultraweak (Sigma) Topology}

\begin{definition}
  \label{def:sigma_top}
  \lean{Ultraweak}
  \leanok
  \uses{def:wstar_def}
  The weak-$*$ topology $\sigma(M,M_{*})$ on a $W^{*}$--algebra $M$ is called the $\sigma$--topology.
\end{definition}

\begin{lemma}
  \label{lem:wstar_unital}
  \lean{CStarAlgebra.isUnital_of_predual}
  \leanok
  \uses{thm:cstar_unital_iff_extreme_Sak_1_6_1}
  Every $W^{*}$--algebra is unital.
\end{lemma}

\begin{proof}
  \leanok
  The norm-closed unit ball of $M=(M_{*})^{*}$ is compact in the $\sigma$--topology and is a
  nonempty compact convex set. By the Krein--Milman theorem it has an extreme point. Theorem
  \ref{thm:cstar_unital_iff_extreme_Sak_1_6_1} then implies that the underlying non-unital
  $C^{*}$--algebra is unital. This argument does not require completeness of the chosen predual.
\end{proof}

\begin{lemma}
  \label{lem:cut_down_sigma_closed_cts} (Sakai 1.7.6)
  \lean{IsStarProjection.Corner.isClosed_ultraweakRange,IsStarProjection.Corner.continuous_ultraweakCutdownₗ}
  \leanok
  \uses{def:sigma_top}
  Let $p$ be any projection in $M$. Then the subalgebra $pMp$ is $\sigma$--closed and the mapping $x\mapsto pxp$ is $\sigma$--continuous.
\end{lemma}

\begin{proof}
  \leanok
  The positive part of the unit ball of $pMp$ is the order interval $[0,p]$, hence is
  $\sigma$--compact. Decomposing an arbitrary element into four positive contractions shows that
  the unit ball of $pMp$ is $\sigma$--compact, so Krein--Smulian gives the closedness of $pMp$.
  For the cutdown map, Sakai's off-diagonal self-adjoint dilation argument shows that both its range
  and kernel are $\sigma$--closed. Since cutdown is an idempotent contraction preserving the unit
  ball, compactness of that ball and the closed range and kernel imply its $\sigma$--continuity.
\end{proof}

\begin{lemma}
  \label{lem:cut_down_of_wstar}
  \lean{IsStarProjection.Corner.wStarAlgebra}
  \leanok
  \uses{def:sigma_top}
  \uses{lem:cut_down_sigma_closed_cts,lem:sigma_closed_star_subalg_wstar}
  If $p\in M$ is a projection, $pMp$ is also a $W^{*}$--algebra with identity $p$.
\end{lemma}

\begin{proof}
  \leanok
  Use the analytically extensible type synonym \texttt{IsStarProjection.Corner} for $pMp$; its
  multiplicative identity is $p$. Its canonical range map is an explicit linear isometry from this
  corner onto its ambient range submodule. Lemma
  \ref{lem:cut_down_sigma_closed_cts} says that this range is ultraweakly closed, so
  \texttt{WStarAlgebra.of\_isClosed\_submodule} supplies a predual for the corner. All transport is
  through these explicit maps rather than definitional equality of subtype structures.
\end{proof}

\begin{lemma}
  \label{lem:left_mul_proj_right_mul_proj_sig_cts_Sak_1_7_7} (Sak 1.7.7)
  \lean{IsStarProjection.continuous_ultraweakMulLeftₗ,IsStarProjection.continuous_ultraweakMulRightₗ}
  \leanok
  \uses{def:sigma_top,lem:cut_down_sigma_closed_cts}
  If $p$ is any projection of $M$, then the mappings $x\mapsto px$ and
  $x\mapsto xp$ are $\sigma$--continuous.
\end{lemma}

\begin{proof}
  \leanok
  The upper off-diagonal contractions $px(1-p)$ obtained from self-adjoint contractions form a
  $\sigma$--compact set. The only nonformal step is to show that this set is closed: a hypothetical
  lower off-diagonal component in a limit contradicts the max-norm identity for orthogonal Peirce
  components after adding a sufficiently large multiple of that component. The unit ball in the
  range of $x\mapsto px$ is then contained in the sum of the compact unit balls of $pMp$ and the
  upper off-diagonal space, and the reverse inclusion holds after intersecting with the ambient unit
  ball. Krein--Smulian closes the range; the kernel is the corresponding range for $1-p$, so the
  idempotent compact-ball criterion gives continuity. Finally, $x\mapsto xp$ is obtained from left
  multiplication by conjugating with the $\sigma$--continuous star operation.
\end{proof}

\begin{lemma}
  \label{lem:star_left_mul_right_mul_sig_cts_Sak_1_7_8} (Sak 1.7.8)
  \lean{Ultraweak.instContinuousStarComplex,Ultraweak.continuous_mulLeftₗ,Ultraweak.continuous_mulRightₗ}
  \leanok
  \uses{def:sigma_top,lem:left_mul_proj_right_mul_proj_sig_cts_Sak_1_7_7,lem:spec_masa_wstar_stonean_Sak_1_7_5,lem:fin_lin_approx_of_stonean_Sak_1_3_1}
  The mappings $x\mapsto x^{*},ax$, and $xa$ are $\sigma$--continuous for
  $x,a \in M$.
\end{lemma}

\begin{proof}
  \leanok
  The self-adjoint and skew-adjoint parts are closed complementary real subspaces, which makes the
  star operation $\sigma$--continuous. More generally, the elements $a$ for which left
  multiplication by $a$ is $\sigma$--continuous form a norm-closed complex submodule: norm
  convergence of the multipliers gives uniform convergence after evaluation by any predual
  functional on the closed unit ball, and Krein--Smulian promotes continuity there to continuity on
  all of $M$. Every projection belongs to this submodule by the preceding lemma. A self-adjoint
  element lies in a masa; its Stonean Gelfand spectrum shows that it is a norm limit of real linear
  combinations of projections. Hence every self-adjoint element, and then every element by real
  and imaginary decomposition, is a continuous left multiplier. Right multiplication is obtained
  by conjugating left multiplication with the continuous star operation.
\end{proof}

\begin{corollary}
  \label{cor:sig_clos_of_cstar_is_wstar_Sak_1_7_9} (Sak 1.7.9)
  \lean{StarSubalgebra.ultraweakClosure,StarSubalgebra.ultraweakClosure_minimal,StarSubalgebra.ultraweakAdjoin,StarSubalgebra.subset_ultraweakAdjoin,StarSubalgebra.isUltraweakClosed_ultraweakAdjoin,StarSubalgebra.ultraweakAdjoin_le,StarSubalgebra.IsUltraweakClosed.wStarAlgebra}
  \leanok
  \uses{lem:left_mul_proj_right_mul_proj_sig_cts_Sak_1_7_7,lem:star_left_mul_right_mul_sig_cts_Sak_1_7_8,def:sigma_top}
  Let $H$ be a subset of $M$ and $\mathcal{A}$ be the $C^{*}$--subalgebra of $M$ generated
  by $H$, and let $\overline{\mathcal{A}}$ be the $\sigma(M,M_{*})$--closure of $\mathcal{A}$. Then
  $\overline{\mathcal{A}}$ is a $W^{*}$--subalgebra of $M$. $\overline{\mathcal{A}}$ is called a $W^{*}$--subalgebra
  of $M$ generated by $H$. Moreover, $\overline{\mathcal{A}}$ is commutative if $\mathcal{A}$ is commutative.
\end{corollary}

\begin{proof}
  \leanok
  The construction \texttt{ultraweakAdjoin} first adjoins $H$ as an ordinary star subalgebra and
  then takes its ultraweak closure. The results \texttt{subset\_ultraweakAdjoin} and
  \texttt{ultraweakAdjoin\_le} characterize it as the least ultraweakly closed star subalgebra
  containing $H$; ultraweak closure itself satisfies the corresponding minimality theorem.
  Ultraweak closedness implies norm closedness and supplies the quotient predual, hence the
  induced C$^*$-algebra is a W$^*$-algebra. Mathlib's closure theorem for semitopological rings
  shows that commutativity is preserved.
\end{proof}

\begin{lemma}
  \label{lem:wstar_of_masa}
  \lean{StarSubalgebra.IsMasa.isUltraweakClosed,StarSubalgebra.IsUltraweakClosed.wStarAlgebra}
  \leanok
  \uses{cor:sig_clos_of_cstar_is_wstar_Sak_1_7_9}
  A maximal abelian $*$--subalgebra $\mathcal{A}$ of a $W^{*}$--algebra $M$ is also a $W^{*}$--algebra.
\end{lemma}

\begin{proof}
  \leanok
  The ultraweak closure of a masa is commutative and contains the masa, so maximality identifies
  it with the masa itself. Thus the masa is ultraweakly closed, and the preceding construction
  gives it a W$^*$-algebra structure.
\end{proof}

\section{Other Topologies on WStarAlgebras}

In this section, let $M$ be a $W^{*}$--algebra.

\begin{proposition}
  \label{prop:loc_cvx_result}
  \lean{Ultraweak.predualDualEquiv}
  \leanok
  The set of $\sigma(M,M_{*})$--continuous linear functionals is precisely $M_{*}$.
\end{proposition}

\begin{proof}
  \leanok
  The explicit linear equivalence \texttt{Ultraweak.predualDualEquiv} sends the specified predual
  onto the continuous dual of the transported ultraweak topology, with evaluation given by the
  original dual pairing.
\end{proof}

\begin{definition}
  \label{def:Mackey_top}
  \lean{MackeySpace,continuous_ofMackey}
  \leanok
  The locally convex topology on $M$ generated by the seminorms $p_{K}(x)=\sup_{\varphi\in K}|\varphi(x)|$, ranging over
  the absolutely convex $\sigma(M_{*},M)$--compact subsets $K$ of $M_{*}$, is called the \textbf{Mackey topology} on $M$.
  We denote this topology by $\tau(M,M_{*})$.
\end{definition}

The Mackey topology is the strongest compatible locally convex topology: its continuous dual is
precisely $M_{*}$, and every other locally convex topology with that continuous dual is weaker.

\begin{definition}
  \label{def:strong_top}
  \lean{Ultraweak.Strong,Ultraweak.Strong.seminorm_apply}
  \leanok
  The locally convex topology on $M$ generated by the seminorms $x\mapsto \varphi(x^{*}x)^{1/2}$
  ranging over the positive $\sigma(M,M_{*})$--continuous functionals $\varphi\in M_{*}$ is
  called the \textbf{strong topology}, or $s$--topology, on $M$. We denote it by $s(M,M_{*})$.
\end{definition}

The second comparison below uses the Mackey--Arens maximality theorem, formalized by
\texttt{continuous\_ofMackey}.
By $A\le B$ on topologies below, we mean $A$ is weaker than $B$.

\begin{theorem}
  \label{thm:sigma_le_s_le_tau}
  \lean{Ultraweak.Strong.continuous_toUltraweakₗ,Ultraweak.Strong.continuous_fromMackeyₗ}
  \leanok
  \uses{def:Mackey_top,def:strong_top,def:sigma_top, prop:extr_sa_sa_unitaries_Sak_1_6_3}
  We have $\sigma(M,M_{*})\le s(M,M_{*})\le \tau(M,M_{*})$.
\end{theorem}

\begin{proof}
  \leanok
  The first comparison is represented by the explicit continuous linear map from the strong
  synonym to the ultraweak synonym.  For the second, the GNS representation shows that every
  strongly continuous functional belongs to the specified predual, so the strong dual is exactly
  the ultraweak dual.  The Mackey--Arens maximality theorem then makes the explicit map from the
  Mackey synonym to the strong synonym continuous.
\end{proof}

\begin{lemma}
  \label{lem:bdd_sigma_cts_iff_bdd_s_cts_Sak_1_8_10} (Sak 1.8.10)
  \lean{Ultraweak.Strong.continuous_ultraweak_iff_strong,Ultraweak.Strong.continuousOn_closedUnitBall_tfae}
  \leanok
  \uses{thm:sigma_le_s_le_tau}
  For a complex-linear functional $\rho$ on $M$, the following are equivalent: $\rho$ is
  $\sigma$--continuous; $\rho$ is $s$--continuous; the restriction of $\rho$ to the norm-closed
  unit ball is $\sigma$--continuous; and that restriction is $s$--continuous.
\end{lemma}

\begin{proof}
  \leanok
  This is an equivalence of continuous linear functionals for the two restricted topologies, not
  an equality of the restricted topologies.  Full continuity follows from continuity on the
  ultraweak unit ball by Krein--Šmulian.  For the strong unit ball, equality of the two continuous
  duals transports closure of the convex intersection of the kernel with that ball to the
  ultraweak topology, where Krein--Šmulian applies again.  In particular, bounded ultraweak
  convergence need not imply strong convergence.
\end{proof}
